% ============================================================================
%  Preambolo — Machine Learning, note riorganizzate
%  Motore: XeTeX (tectonic). Font: XCharter (testo e formule), Fira Sans (titoli).
% ============================================================================

% --- Font e lingua -----------------------------------------------------------
\usepackage{amsmath,amsthm}
\usepackage{fontspec}
\usepackage{unicode-math}
\setmainfont{XCharter}[Extension=.otf, UprightFont=*-Roman, BoldFont=*-Bold,
  ItalicFont=*-Italic, BoldItalicFont=*-BoldItalic]
\setsansfont{FiraSans}[Extension=.otf, UprightFont=*-Regular, BoldFont=*-SemiBold,
  ItalicFont=*-Italic, BoldItalicFont=*-SemiBoldItalic, Scale=MatchLowercase]
\setmathfont{XCharter-Math.otf}
\usepackage{polyglossia}
\setmainlanguage{italian}
\usepackage{microtype}
\linespread{1.06}

% --- Pagina ------------------------------------------------------------------
\usepackage[a4paper, left=2.4cm, right=2.4cm, top=2.5cm, bottom=2.6cm,
  headheight=15pt, headsep=0.6cm, footskip=1cm]{geometry}
\setlength{\parindent}{0pt}
\setlength{\parskip}{0.55ex plus 0.3ex minus 0.1ex}
\setlength{\emergencystretch}{2.5em}
\raggedbottom
\clubpenalty=8000 \widowpenalty=8000

% --- Colori (usati solo dove serve attirare l'attenzione e nei grafici) ------
\usepackage{xcolor}
\definecolor{accento}{HTML}{A3262A}   % avvertenze
\definecolor{inchiostro}{HTML}{1A1A1A}
\definecolor{grigio}{HTML}{6E6E6E}
\definecolor{grigiochiaro}{HTML}{EFEFEF}
\definecolor{blu}{HTML}{1F5A9E}
\definecolor{rosso}{HTML}{C0392B}
\definecolor{verde}{HTML}{2E7D32}
\definecolor{arancio}{HTML}{D9822B}
\definecolor{viola}{HTML}{6A3D9A}

% --- Titoli ------------------------------------------------------------------
\usepackage{titlesec}
\titleformat{\part}[display]{\centering\sffamily}
  {\large\scshape\color{grigio}Parte \thepart}{1.2ex}{\Huge\bfseries}
\titleformat{\chapter}[display]{\sffamily}
  {\large\color{grigio}Capitolo \thechapter}{0.6ex}
  {\huge\bfseries}[\vspace{0.8ex}{\color{inchiostro}\titlerule[0.8pt]}]
\titlespacing*{\chapter}{0pt}{-1.2cm}{1.4cm}
\titleformat{\section}{\sffamily\Large\bfseries}{\thesection}{0.8em}{}
\titlespacing*{\section}{0pt}{3.2ex plus 1ex minus .4ex}{1.4ex plus .3ex}
\titleformat{\subsection}{\sffamily\large\bfseries}{\thesubsection}{0.8em}{}
\titlespacing*{\subsection}{0pt}{2.6ex plus .8ex minus .3ex}{1ex plus .2ex}
\titleformat{\subsubsection}[runin]{\sffamily\bfseries}{}{0em}{}[.]
\titlespacing*{\subsubsection}{0pt}{1.6ex plus .5ex minus .2ex}{0.6em}
\setcounter{secnumdepth}{2}
\setcounter{tocdepth}{1}

% --- Intestazioni ------------------------------------------------------------
\usepackage{fancyhdr}
\pagestyle{fancy}
\fancyhf{}
\makeatletter
\renewcommand{\chaptermark}[1]{\markboth{\ifnum\c@chapter>0 \thechapter.\ \fi#1}{}}
\makeatother
\renewcommand{\sectionmark}[1]{\markright{\thesection\ \ #1}}
\fancyhead[L]{\sffamily\small\color{grigio}\leftmark}
\fancyhead[R]{\sffamily\small\color{grigio}\rightmark}
\fancyfoot[C]{\sffamily\small\thepage}
\renewcommand{\headrulewidth}{0.4pt}
\fancypagestyle{plain}{\fancyhf{}\fancyfoot[C]{\sffamily\small\thepage}\renewcommand{\headrulewidth}{0pt}}

% --- Liste, tabelle, didascalie ----------------------------------------------
\usepackage{enumitem}
\setlist{itemsep=0.25ex, topsep=0.6ex, parsep=0.2ex, leftmargin=1.6em}
\setlist[enumerate]{label=\arabic*.}
\usepackage{booktabs, array, tabularx, multirow}
\renewcommand{\arraystretch}{1.18}
\usepackage[font=small, labelfont={bf,sf}, labelsep=period, margin=0.6cm]{caption}
\captionsetup[table]{position=top, skip=4pt}
\usepackage[normalem]{ulem}

% --- Grafica -----------------------------------------------------------------
\usepackage{tikz}
\usetikzlibrary{arrows.meta, positioning, calc, decorations.pathreplacing, patterns,
  shapes.geometric, shapes.misc, fit, backgrounds, matrix, 3d, intersections}
\usepackage{pgfplots}
\pgfplotsset{compat=1.18}
\usepgfplotslibrary{fillbetween, groupplots}
\pgfplotsset{
  stile/.style={
    axis lines=left, tick style={grigio}, axis line style={-{Stealth[length=2mm]}, inchiostro},
    every tick label/.append style={font=\footnotesize},
    label style={font=\small}, legend style={font=\footnotesize, draw=none, fill=none},
    legend cell align=left, clip mode=individual, samples=150,
  },
  scatola/.style={
    axis lines=box, tick style={grigio}, every tick label/.append style={font=\footnotesize},
    label style={font=\small}, legend style={font=\footnotesize, draw=grigio!50, fill=white},
    legend cell align=left, samples=150,
  },
  nudo/.style={axis lines=none, ticks=none},
  empty line=jump,
}
\tikzset{
  >={Stealth[length=2.2mm]},
  blocco/.style={draw, rounded corners=1.5pt, minimum height=8mm, inner sep=4pt,
    align=center, font=\small},
  neur/.style={circle, draw, minimum size=7mm, inner sep=0pt, font=\small},
  neurp/.style={circle, draw, fill=grigiochiaro, minimum size=7mm, inner sep=0pt, font=\small},
  freccia/.style={->, thick},
  cpos/.style={circle, fill=blu, inner sep=1.4pt},
  cneg/.style={circle, draw=rosso, fill=white, thick, inner sep=1.4pt},
}

% --- Riquadri ----------------------------------------------------------------
\usepackage[most]{tcolorbox}
\tcbset{
  base/.style={enhanced, breakable, sharp corners, boxrule=0pt, frame hidden,
    left=8pt, right=6pt, top=3pt, bottom=3pt, before skip=1.8ex plus .5ex,
    after skip=1.8ex plus .5ex, fontupper=\normalsize},
}
\newcommand{\etichetta}[2]{{\sffamily\bfseries\color{#1}#2}}
% Definizione: filetto nero a sinistra
\newtcolorbox{definizione}[1][]{base, colback=white,
  borderline west={1.6pt}{0pt}{inchiostro},
  before upper={\etichetta{inchiostro}{Definizione\ifx\relax#1\relax\else\ --- #1\fi.}\ }}
% Teorema: cornice sottile su tutti i lati
\newtcolorbox{teorema}[1][]{base, colback=white, frame hidden,
  borderline={0.6pt}{0pt}{inchiostro},
  left=7pt, right=7pt, top=5pt, bottom=5pt,
  before upper={\etichetta{inchiostro}{Teorema\ifx\relax#1\relax\else\ --- #1\fi.}\ \itshape}}
\newtcolorbox{proprieta}[1]{base, colback=white, frame hidden,
  borderline={0.6pt}{0pt}{inchiostro},
  left=7pt, right=7pt, top=5pt, bottom=5pt,
  before upper={\etichetta{inchiostro}{#1.}\ }}
% Esempio: filetto grigio sottile
\newtcolorbox{esempio}[1][]{base, colback=white,
  borderline west={0.8pt}{0pt}{grigio},
  before upper={\etichetta{grigio}{Esempio\ifx\relax#1\relax\else\ --- #1\fi.}\ }}
% Attenzione: l'unico riquadro colorato
\newtcolorbox{attenzione}[1][]{base, colback=accento!4,
  borderline west={1.6pt}{0pt}{accento},
  before upper={\etichetta{accento}{Attenzione\ifx\relax#1\relax\else\ --- #1\fi.}\ }}
% Algoritmo: fondo grigio chiaro
\newtcolorbox{algoritmo}[1][]{base, colback=grigiochiaro,
  before upper={\etichetta{inchiostro}{Algoritmo\ifx\relax#1\relax\else\ --- #1\fi.}\par\smallskip}}
% Idea / osservazione chiave
\newtcolorbox{idea}[1][]{base, colback=white,
  borderline west={0.8pt}{0pt}{inchiostro}, borderline west={0.8pt}{2.4pt}{inchiostro},
  left=11pt,
  before upper={\etichetta{inchiostro}{\ifx\relax#1\relax Idea chiave\else #1\fi.}\ }}
% Apertura e chiusura dei capitoli
\newtcolorbox{inbreve}{base, colback=white, left=0pt, right=0pt, top=0pt, bottom=4pt,
  fontupper=\itshape\normalsize}
\newtcolorbox{puntichiave}{base, colback=grigiochiaro, top=5pt, bottom=5pt,
  before upper={\etichetta{inchiostro}{Punti chiave}\par\smallskip}}
\newtcolorbox{domande}{base, colback=white, borderline={0.4pt}{0pt}{grigio}, top=5pt, bottom=5pt,
  before upper={\etichetta{inchiostro}{Domande di verifica}\par\smallskip}}

% Dimostrazioni
\renewcommand{\qedsymbol}{\rule{1.2ex}{1.2ex}}
\makeatletter
\renewenvironment{proof}[1][\proofname]{\par\pushQED{\qed}\normalfont
  \topsep6\p@\@plus6\p@\relax\trivlist
  \item[\hskip\labelsep{\sffamily\bfseries #1.}]\ignorespaces}{\popQED\endtrivlist\@endpefalse}
\makeatother
\AtBeginDocument{\renewcommand{\proofname}{Dimostrazione}\renewcommand{\figurename}{Figura}\renewcommand{\tablename}{Tabella}}

% --- Comandi matematici ------------------------------------------------------
\newcommand{\x}{\mathbf{x}}
\newcommand{\w}{\mathbf{w}}
\newcommand{\y}{\mathbf{y}}
\newcommand{\dd}{\mathbf{d}}
\newcommand{\h}{\mathbf{h}}
\newcommand{\vv}[1]{\mathbf{#1}}
\newcommand{\R}{\mathbb{R}}
\newcommand{\T}{^{T}}
\newcommand{\norm}[1]{\lVert #1\rVert}
\newcommand{\abs}[1]{\lvert #1\rvert}
\newcommand{\E}{\mathbb{E}}
\DeclareMathOperator*{\argmin}{arg\,min}
\DeclareMathOperator*{\argmax}{arg\,max}
\DeclareMathOperator{\sign}{sign}
\DeclareMathOperator{\Var}{Var}
\DeclareMathOperator{\Cov}{Cov}
\DeclareMathOperator{\Bias}{Bias}
\newcommand{\Remp}{R_{\mathrm{emp}}}
\newcommand{\VC}{\mathrm{VC}}
\newcommand{\net}{\mathit{net}}
\newcommand{\out}{\mathit{out}}
\newcommand{\balpha}{\symbf{\alpha}}
\newcommand{\bxi}{\symbf{\xi}}
\newcommand{\bphi}{\symbf{\Phi}}
\newcommand{\eps}{\varepsilon}

% Testo
\newcommand{\termine}[1]{\textbf{#1}}       % termine definito per la prima volta
\newcommand{\ingl}[1]{\textit{#1}}          % termine inglese
\newcommand{\rif}[1]{\S\,\ref{#1}}
\newcommand{\figrif}[1]{Figura~\ref{#1}}

% --- Collegamenti ------------------------------------------------------------
\usepackage[hidelinks, bookmarksnumbered, pdfusetitle]{hyperref}
\usepackage{bookmark}
